\documentclass[11pt]{article}
\usepackage[utf8]{inputenc}
\usepackage[T1]{fontenc}
\usepackage[margin=1in]{geometry}
\usepackage[hidelinks]{hyperref}
\usepackage{parskip}
\pagestyle{empty}

\begin{document}

\noindent\textbf{\large Cover Letter --- \textit{npj Complexity}}

\vspace{1em}

\noindent Dear Editors,

\vspace{0.6em}

I am pleased to submit our manuscript, \textbf{``An Area-Preserving H\'enon-Map Model for the Riemann Zeros: A Deterministic-Dynamics Approach with Quantum and Dissipative Solvers,''} for consideration as a research article in \textit{npj Complexity}.

\vspace{0.4em}

\noindent\textbf{What the paper does.} We ask a question at the interface of nonlinear dynamics, complex systems, and number theory: can a single low-dimensional \emph{deterministic} dynamical system, without hand-tuned unfolding, reproduce both the global counting function of the Riemann zeros and their local level repulsion? The study is a direct follow-up to our recently \textbf{published} result establishing that the symbolic dynamics of the Logistic map at its band-merging point is isomorphic to the prime sieve, with the Hardy--Littlewood twin-prime constant emerging as a dynamical fixed point (Wang, \textit{Research in Mathematics} 13(1), 2026). Given that peer-reviewed foundation, we take the natural Hilbert--P\'olya step and analyze the \emph{spectrum} of the associated map. Because the 1D Logistic map is dissipative and cannot host a unitary spectrum, we lift it to the 2D area-preserving H\'enon map and study it near the edge of chaos with two numerically independent eigensolvers (a unitary Fourier solver and a Markovian dissipative solver).

\vspace{0.4em}

\noindent\textbf{Why \textit{npj Complexity}.} The work sits squarely in the journal's scope --- emergence, self-organization, and nonlinear/complex dynamics --- and uses the complex-systems toolkit (edge-of-chaos dynamics, KAM breakup, spectral statistics, large-scale computation) to address a classical problem from a fresh, constructive angle. It also exemplifies a data-driven, ``bottom-up'' discovery workflow that we believe will interest the journal's readership.

\vspace{0.4em}

\noindent\textbf{On rigor and claims.} We are deliberately careful about epistemic status. The manuscript presents its results as \emph{numerical observations and heuristic arguments, not proofs}, and includes an explicit table classifying every claim as established, published, numerical, heuristic, qualitative, or conjectural. We do not claim to have identified the Hilbert--P\'olya operator or to have proven any statement about the Riemann Hypothesis. All code, data, and logs are openly available (\url{https://github.com/maris205/riemann_henon}) for full reproducibility.

\vspace{0.4em}

\noindent\textbf{Originality and prior consideration.} The manuscript is original, has not been published elsewhere, and is not under consideration by another journal. The sole author has no competing interests.

\vspace{0.6em}

\noindent Thank you for considering our submission. We would be happy to suggest suitable referees with expertise in quantum chaos, random matrix theory, and dynamical systems upon request.

\vspace{1em}

\noindent Sincerely,

\vspace{0.3em}

\noindent Liang Wang \\
School of Artificial Intelligence and Automation, \\
Huazhong University of Science and Technology, Wuhan, P.R. China \\
\href{mailto:wangliang.f@gmail.com}{wangliang.f@gmail.com}

\end{document}
